\chapter{Conclusions and Future Work}
\label{chap:conclusiones}

% =============================================================================
\section{General Conclusions}
\label{sec:conclusions-general}
% =============================================================================

This thesis has investigated a fundamental tension at the intersection of
computational efficiency and representational fidelity in attention-based
building change detection from very high-resolution satellite imagery: whether
the quadratic memory bottleneck of standard Multi-Head Self-Attention can be
eliminated through a principled rank-augmentation of its key-value buffer,
and whether the resulting linear-complexity mechanism can be substituted
drop-in within a state-of-the-art hierarchical Siamese Transformer without
incurring meaningful accuracy degradation across a diverse multi-dataset
evaluation suite. The experimental program documented in
Chapter~\ref{chap:experiments} provides a comprehensive and affirmative answer
to both questions.

The central architectural contribution is RALA-ChangeFormer: a hierarchical
Siamese encoder-decoder for bi-temporal building change detection in which every
Multi-Head Self-Attention block within the SegFormer-B2 encoder is replaced by
a Rank-Augmented Linear Attention block \cite{Fan2025RALA}. The substitution
maps the attention-related time complexity from $\mathcal{O}(N_l^2 d)$ to
$\mathcal{O}(N_l d^2)$ and the memory complexity from $\mathcal{O}(N_l^2)$ to
$\mathcal{O}(d^2)$ at each encoder level $l$, where $N_l$ is the number of
spatial tokens and $d$ is the per-head embedding dimension. Because $N_l \gg d$
for all four encoder levels when processing VHR patches at $256 \times 256$
pixels and above, this asymptotic reduction is practically significant: it
eliminates the $N_l \times N_l$ attention map from GPU memory during both the
forward and backward passes, replacing it with the $d \times d$ rank-augmented
key-value buffer $B$. The theoretical foundation of this substitution, developed
in Section~\ref{sec:rala-theory}, establishes that rank collapse of the buffer---
the principal failure mode of standard linear attention in dense prediction
settings---is suppressed through two complementary mechanisms: the global query
descriptor $Q_g$, which introduces token-relevance weighting that differentially
amplifies informationally distinct outer products into the buffer accumulation,
and the Hadamard output modulation, which injects local per-token information
into the output feature representation and breaks the rank inequality that would
otherwise collapse the output matrix rank to the minimum of $\mathrm{rank}(\kappa(Q))$
and $\mathrm{rank}(B)$.

The primary empirical validation of this thesis is the demonstration of
statistical accuracy parity between RALA-ChangeFormer and the full quadratic
ChangeFormer \cite{Bandara2022ChangeFormer} across all six evaluated benchmarks.
Across the complete evaluation matrix---encompassing LEVIR-CD, DSIFN-CD, WHU-CD,
LEVIR-CD+, SYSU-CD, and S2Looking---the mean F1 advantage of RALA-ChangeFormer
over ChangeFormer is $+0.14$ percentage points, a margin smaller than the
run-to-run variance of a single training configuration and devoid of statistical
significance. This result holds consistently across nadir-view and oblique-view
acquisition geometries, binary and multi-class change annotation schemes, and
spatial resolutions spanning two orders of magnitude from $0.3\,\mathrm{m}$ to
$0.5\,\mathrm{m}$ per pixel. The accuracy parity is confirmed at both the summary
metric level and the per-error-type qualitative level: the error map distributions
of RALA-ChangeFormer exhibit the same morphological profiles as ChangeFormer on
each dataset, without introducing any failure category specific to the
rank-augmented attention mechanism.

The computational efficiency gains are equally decisive. As established through
controlled hardware profiling at three input patch resolutions, RALA-ChangeFormer
delivers a \textbf{21.8\% reduction in peak GPU VRAM consumption} relative to
ChangeFormer (10.38~GB versus 13.28~GB at $256 \times 256$, batch size 16) and
an \textbf{8.0\% reduction in inference GFLOPs} (185.93 versus 202.17~GFLOPs
per image pair). These reductions are accompanied by a \textbf{50\% expansion
in the maximum feasible training batch size} before Out-Of-Memory failure, a
ratio that is maintained uniformly across the tested resolution range from
$256 \times 256$ to $1024 \times 1024$ pixels. In practical terms, ChangeFormer
at standard operating resolution requires a minimum of 14~GB GPU VRAM and
restricts training at the native $1024 \times 1024$ LEVIR-CD resolution to a
batch size of 2; RALA-ChangeFormer operates within 12~GB GPU VRAM and expands
the $1024 \times 1024$ batch limit to 3. These numbers translate directly into
a lower hardware entry threshold for Transformer-based BCD research: institutions
operating with consumer-grade GPUs in the 10--16~GB VRAM range---representative
of the computing infrastructure available at the majority of academic Earth
observation groups in Latin America, including those engaged in PeruSAT-1 urban
monitoring---can now conduct full-resolution Siamese Transformer training within
a single GPU budget that previously supported only convolutional or
token-compressed architectures.

Taken together, these findings confirm that rank-augmented linear attention
constitutes a principled, non-heuristic pathway for democratizing access to
state-of-the-art VHR building change detection. The efficiency gains are not
obtained through approximations that degrade accuracy under distribution shift,
compression schemes that reduce parameter capacity, or knowledge distillation
from larger models: they are a direct algebraic consequence of replacing the
dense $N \times N$ attention map with a $d \times d$ buffer whose rank is
actively maintained against collapse through the learned weighting and modulation
mechanisms of RALA. This structural clarity is the thesis's most enduring
methodological contribution, as it provides a framework for systematic
efficiency-accuracy analysis that is directly transferable to other hierarchical
dense prediction architectures that share the quadratic attention bottleneck.


% =============================================================================
\section{Limitations of Low-Rank Attention Approximations}
\label{sec:conclusions-limitations}
% =============================================================================

The experimental program also identifies a set of structurally well-characterized
limitations that delineate the operational boundaries of the RALA-ChangeFormer
framework and motivate the future research directions articulated in
Section~\ref{sec:conclusions-future}.

The most severe performance boundary encountered in the six-dataset evaluation
is the degradation on S2Looking, where both ChangeFormer and RALA-ChangeFormer
achieve only 65.55\% and 65.91\% F1 respectively---a gap of approximately
24--25 points relative to LEVIR-CD performance. The qualitative analysis of
Section~\ref{sec:exp-dynamics} establishes that this degradation is not a
consequence of the rank-augmented attention mechanism: the error maps of both
architectures exhibit statistically identical parallax halo artifacts, concentrated
at the perimeters of tall unchanged buildings where the apparent positional shift
between acquisitions from slightly different orbital tracks is largest. This
failure mode is architectural rather than attentional: it originates in the
shared SegFormer-B2 convolutional encoder backbone, which extracts features from
observed pixel intensity patterns without access to the satellite viewing angle
metadata, off-nadir angle pairs, or stereophotogrammetric depth estimates that
would be necessary to normalize out the apparent positional offsets induced by
the oblique acquisition geometry. No modification to the attention mechanism
alone---whether quadratic or linear, whether with or without rank augmentation---
can resolve this geometric inconsistency, because the encoder's input tokens
themselves are contaminated by parallax-induced apparent displacement before
any attention computation occurs. The implication is clear: effective BCD on
side-looking satellite imagery requires geometric pre-processing stages that
are currently absent from all standard training pipelines, including that of
this thesis.

A second category of limitation is exposed by the SYSU-CD results. The dataset's
multi-class annotation scheme assigns positive change labels to road alterations,
vegetation removal, sea reclamation, and building construction simultaneously,
while the model is trained under a binary BCE-Dice loss that was calibrated on
the building-construction-dominated training distribution of LEVIR-CD. The result
is a systematic false alarm elevation in regions containing semantically valid
but non-building change events---visible as dense red patches in the SYSU-CD
error maps---that neither ChangeFormer nor RALA-ChangeFormer can suppress without
loss of recall on the building class. This limitation is not specific to RALA;
it reflects the fundamental mismatch between a building-centric training objective
and a multi-class evaluation target. Resolving it requires either multi-class
loss supervision during training, domain adaptation from the binary to the
multi-class setting, or dataset-specific fine-tuning that explicitly targets the
SYSU-CD annotation taxonomy. In all these cases, the limitation is a property
of the learning objective design rather than of the attention architecture.


A third category of limitation relates to the encoding capacity of the shared backbone, which integrates convolutional and Transformer layers. As documented in the parameter and FLOPs analysis of Table~\ref{tab:efficiency}, both ChangeFormer and RALA-ChangeFormer operate at 41~M parameters, a fixed capacity determined by the SegFormer-B2 backbone design. This capacity ceiling bounds the upper F1 that either model can achieve across all six datasets, and the performance gradient observed between the two architectures---ranging from $+0.13$ F1 points on LEVIR-CD to $+0.36$ on S2Looking---reflects the degree to which the rank augmentation provides an attention-level representational advantage over the uniform linear baseline within the shared backbone capacity. When the backbone itself is the
binding constraint---as on S2Looking, where the geometric normalization of the
input data is the operative limitation---further improvements to the attention
mechanism can yield only marginal gains. This observation implies that
architectural advances in efficient BCD must target both the attention
computation and the encoder input pipeline simultaneously to unlock the next
tier of accuracy improvements on geometrically challenging benchmarks.

Finally, the practical memory and compute floor of RALA-ChangeFormer, while
substantially lower than that of its quadratic counterpart, remains non-trivial
in absolute terms. At $1024 \times 1024$ native resolution, the maximum feasible
batch size of 3 continues to produce noisy gradient estimates that constrain
convergence speed relative to the batch sizes achievable at $256 \times 256$.
For institutions with access to only 8~GB GPU VRAM---still a common constraint
in the Global South---the architecture remains out of reach for full-resolution
training even with RALA, and patch-based training with reduced resolution must
be used. These constraints reinforce the argument that additional algorithmic
efficiency innovations, beyond the attention mechanism substitution alone, are
needed to make end-to-end full-resolution VHR BCD universally accessible.


% =============================================================================
\section{Future Lines of Research}
\label{sec:conclusions-future}
% =============================================================================

The limitations documented in Section~\ref{sec:conclusions-limitations} point
to a coherent research roadmap organized around two complementary pillars: the
technical extension of the efficient attention framework through adaptive and
geometry-aware mechanisms, and the geographic and morphological expansion of the
evaluation and deployment context toward the informal urban environments that
characterize Latin American metropolitan areas.

\subsection*{Technical Extensions: Adaptive Rank Control and Geometric Alignment}

The most immediate technical limitation of RALA is that the rank augmentation
strength---governed by the global query descriptor $Q_g$ and the token relevance
weights $\alpha_j$---is a fixed learned function of the input data, without any
explicit mechanism for adapting the effective rank of the key-value buffer $B$
to the local information content of the scene. In spectral environments with
high token diversity, such as informal urban areas with heterogeneous rooftop
materials and irregular spatial configurations, the current buffer may still
experience partial rank loss when the query distribution is insufficiently
diverse. A principled response to this is the integration of an online
lightweight singular value decomposition (SVD) tracking procedure into the
attention block: by monitoring the singular value spectrum of $B$ across training
iterations, the model can dynamically identify when the effective rank of the
buffer is falling below a target threshold and modulate the attention mixing
parameters accordingly. This data-dependent rank adaptation heuristic would
extend RALA's representational robustness to scenes that differ systematically
from the token diversity profiles encountered in the standard benchmark training
distributions, without requiring any increase in the $\mathcal{O}(N d^2)$
asymptotic complexity of the attention computation.

The S2Looking results motivate a distinct line of extension: the design of
upstream geometric perspective pre-alignment modules that normalize the apparent
positional offsets induced by oblique satellite acquisition before the encoder
processes the bi-temporal image pair. A promising direction is the integration
of a lightweight stereophotogrammetric depth estimation head---trained on
auxiliary digital surface model data available from global topographic archives
such as SRTM or Copernicus DEM---that estimates per-building height maps from
single-date VHR imagery and uses them to predict and correct the expected
parallax displacement for a given off-nadir viewing angle pair. This geometric
normalization stage would be inserted between the preprocessing pipeline and the
patch embedding module, operating on the raw image pair before tokenization and
producing geometrically corrected bi-temporal patches that are free of the
apparent positional offsets that currently confound the encoder's feature
comparison. Such a pre-alignment module is architecture-agnostic---it would
benefit both quadratic and linear attention variants equally---but its
development is motivated specifically by the parallax degradation documented
for the side-looking benchmark in this thesis.

A further technical direction concerns the extension of RALA to multi-temporal
sequences beyond the bi-temporal pair format. The current formulation processes
two acquisitions through a weight-sharing Siamese encoder and computes pairwise
feature differences at each scale, a design that is sufficient for detecting
change between two epochs but cannot represent the temporal trajectory of change
events across three or more acquisitions. Adapting the global query descriptor
to aggregate information across multiple temporal tokens---rather than the two
temporal branches of the current Siamese design---would enable RALA-ChangeFormer
to model progressive construction sequences, detect demolition followed by
reconstruction, and assess seasonal temporal stability of detected changes within
a unified architectural framework. This extension would require a more general
formulation of the $\alpha_j$ weighting mechanism to account for the asymmetric
temporal relationships between an arbitrary number of acquisitions, and its
memory footprint would remain $\mathcal{O}(N d^2)$ per attention block provided
that the buffer accumulation is performed independently across temporal channels.

\subsection*{Geographic and Morphological Expansion: Benchmarks for the
Global South}

Beyond technical extensions of the RALA mechanism itself, this thesis identifies
a deeper structural challenge in the BCD literature: the systematic absence of
evaluation benchmarks representative of the informal urban growth patterns,
mixed-use land parcel structures, and irregular building topologies that
characterize the metropolitan areas of Latin America, Sub-Saharan Africa, and
South and Southeast Asia---regions that are simultaneously among the most
rapidly urbanizing on Earth and the most underserved by operational BCD systems. The six benchmark datasets used in this thesis share a common acquisition profile
that reflects the research priorities of communities with access to large-scale
VHR archives and annotation capacity: near-nadir viewing geometry, binary or
categorically simple change labels, predominantly planned or semi-planned urban
morphologies in East Asian cities, and spatial resolutions calibrated for
commercial satellite imagery. Lima, Metropolitan Lima---with its simultaneous
horizontal periphery expansion into desert hillsides and vertical densification
of its consolidated residential districts---represents a qualitatively different
change detection challenge from those captured by LEVIR-CD or WHU-CD. Peripheral
hillside settlements in districts such as Villa María del Triunfo or San Juan de
Miraflores grow through irregular plot subdivision, incremental self-construction
across multiple temporal epochs, and building typologies that differ
fundamentally from the rectangular footprint, flat roof, and homogeneous material
signature that characterizes the dominant positive class in existing benchmarks.
These structural differences mean that models trained exclusively on existing
international benchmarks may exhibit systematic failure modes when applied to
Peruvian urban monitoring contexts that are not captured by any available
evaluation protocol.

Addressing this gap requires the creation of a dedicated bi-temporal VHR building
change detection benchmark for Peruvian and, more broadly, Latin American urban
environments, constructed from PeruSAT-1 imagery supplemented by high-resolution
commercial acquisitions over Lima, São Paulo, Bogotá, and other large metropolitan
areas with active informal growth. The benchmark design must engage directly with
the annotation challenges that are specific to this context: the distinction
between informal incremental construction---where a roof extension added above a
ground-floor dwelling constitutes a structurally significant change that would
not be labeled as change in a binary footprint-change dataset---and the permanent
replacement of one structure by another; the co-occurrence of formal and informal
construction within a single image patch, requiring a nuanced annotation ontology
that existing binary labeling conventions cannot represent; and the calibration
of evaluation metrics to reflect the planning and governance priorities of the
institutions that would consume such a system's outputs, which differ
substantially from the precision-recall trade-offs that drive optimization on
academic benchmarks. The GINIA urban data analysis subgroup at UTEC is uniquely
positioned to lead this initiative, combining access to PeruSAT-1 imagery,
domain expertise in Peruvian urban dynamics, and the methodological foundation
developed in this thesis.

At the operational deployment level, future work must address the integration
of RALA-ChangeFormer into a complete urban monitoring workflow: geometric and
radiometric preprocessing of PeruSAT-1 imagery at continental scale, inference
at the level of entire city districts rather than $256 \times 256$ patches,
mosaicking of patch-level predictions into spatially consistent city-wide change
maps, and the generation of cartographic products suitable for consumption by
public institutions engaged in land use planning, informal settlement monitoring,
and seismic vulnerability assessment. Model compression techniques---including
quantization-aware training, structured pruning, and knowledge distillation into
compact convolutional student models---would further reduce the inference memory
footprint to levels compatible with cloud-native or edge deployment architectures,
extending the operational reach of the RALA-ChangeFormer framework to
environments where neither A100 GPU clusters nor constant connectivity can be
assumed.

The convergence of these technical and geographic research directions points toward a research programme of sustained ambition: the construction of computationally accessible, geographically representative, and operationally deployable deep learning systems for bi-temporal urban change monitoring across the full morphological diversity of the Global South. This thesis establishes the efficient attention mechanism foundations on which such a programme can be built.
